\section{Erkinlik darajalari}


\begin{enumerate}
\item Bir atomli ideal gazning erkinlik darajasi nechaga teng?
\begin{enumerate}[label=\Alph*)]
\item 2
\item 3
\item 5
\item 6
\end{enumerate}

\item Ikki atomli gaz uchun (tebranma harakat hisobga olinmasa) erkinlik darajasi nechaga teng?
\begin{enumerate}[label=\Alph*)]
\item 3
\item 4
\item 5
\item 6
\end{enumerate}

\item Uch atomli ideal gazning erkinlik darajasi nechaga teng?
\begin{enumerate}[label=\Alph*)]
\item 3
\item 4
\item 5
\item 6
\end{enumerate}

\item 2 mol bir atomli gaz 300 K temperaturada turibdi. Ichki energiyasini toping (kJ).
\begin{enumerate}[label=\Alph*)]
\item 5.0
\item 7.5
\item 10.0
\item 12.5
\end{enumerate}

\item 1 mol ikki atomli gazning temperaturasi 200 K ga oshirilsa, ichki energiyasi qancha ortadi (kJ)?
\begin{enumerate}[label=\Alph*)]
\item 2.08
\item 4.16
\item 6.24
\item 8.32
\end{enumerate}

\item Gazning ichki energiyasi ($U=\frac{i}{2}nRT$) formulaga ko‘ra aniqlanadi. Bu yerda (i) nimani bildiradi?
\begin{enumerate}[label=\Alph*)]
\item Bosimni
\item Hajmni
\item Erkinlik darajasini
\item Zichlikni
\end{enumerate}

\item 3 mol uch atomli gaz 400 K temperaturada turibdi. Ichki energiyasini toping (kJ).
\begin{enumerate}[label=\Alph*)]
\item 14.9
\item 29.9
\item 44.9
\item 59.9
\end{enumerate}

\item Bir atomli va ikki atomli gazlarning mollari hamda temperaturasi bir xil. Ichki energiyalari nisbati qanday?
\begin{enumerate}[label=\Alph*)]
\item 3:5
\item 5:3
\item 2:1
\item 1:2
\end{enumerate}

\item 2 mol bir atomli gazning temperaturasi 150 K ga oshirilganda ichki energiyasi qancha ortadi (kJ)?
\begin{enumerate}[label=\Alph*)]
\item 1.87
\item 3.74
\item 5.61
\item 7.48
\end{enumerate}

\item Ideal gazning ichki energiyasi 18.7 kJ, modda miqdori 2.5 mol va temperaturasi 300 K. Gazning erkinlik darajasini aniqlang.
\begin{enumerate}[label=\Alph*)]
\item 3
\item 4
\item 5
\item 6
\end{enumerate}
\end{enumerate}



\textbf{Javoblar kaliti}

\begin{center}
\begin{tabular}{|c|c|c|c|c|c|c|c|c|c|}
\hline
1 & 2 & 3 & 4 & 5 & 6 & 7 & 8 & 9 & 10 \\
\hline
B & C & D & B & B & C & D & A & B & D \\
\hline
\end{tabular}
\end{center}